\chapter{Honest Limitations}

This chapter exists because a backtester that only shows its wins is the most
common way a systematic-trading project deceives its own author. Read it before
reading Chapter~8 as encouragement rather than as a demonstration of technique.

\section{Everything in Chapter 8 is synthetic}

The demo's Sharpe of 0.62 is measured against data generated by this project's
own \texttt{SyntheticSource}, tuned (Chapter~2) so the momentum signal has
something real to find. It says the plumbing is correct -- data flows through
features, signals, validation, costs, and risk without leaking the future or
silently dropping days. It says \emph{nothing} about whether momentum, or any
other signal in this repository, would make money on real markets. No backtest
on synthetic data can say that, by construction.

\section{Retail systematic trading is a near-zero-sum game}

Even against real data, an individual running this pipeline is competing
against firms with colocated servers, direct exchange feeds, proprietary
alternative data, and teams of specialists across the four roles described in
the companion research note (\texttt{quant-trading-reality.md}). None of that
can be replicated by adding more code to this repository. The honest reading
of ``can I build my own code and profit from it'' is: not reliably, and the
expected value of trying, net of the hundreds of hours it takes, is close to
zero to negative for most people, most of the time.

\section{Why the LSTM (and the GBR) should not be trusted here}

\begin{itemize}
  \item \textbf{Low signal-to-noise.} Even the synthetic alpha, deliberately
  exaggerated, produces a modest Sharpe once realistic costs are charged. Real
  single-name equity alphas are far weaker and far more crowded.
  \item \textbf{Non-stationarity.} A model trained on one regime (Chapter~2's
  bull/bear/chop states) has no guarantee its learned relationship survives
  into the next regime -- and real markets change regime on no fixed schedule.
  \item \textbf{Multiple-comparisons risk.} Trying enough feature sets,
  architectures, and hyperparameters against the same historical window will
  eventually produce an attractive-looking backtest by chance alone. Walk-forward
  discipline (Chapter~5) reduces this but does not eliminate it if the model
  itself, or its hyperparameters, were chosen by looking at walk-forward
  results across many attempts.
  \item \textbf{Costs are not an afterthought.} Chapter~8 shows the ML variant
  had \emph{higher} turnover (0.17 vs.\ 0.10) as well as worse gross performance
  -- a common pattern: models that chase noisy short-horizon signals trade more
  and pay more in the process.
\end{itemize}

\section{What this project is actually worth}

\begin{enumerate}
  \item \textbf{A portfolio artefact.} A clean backtester, an honest walk-forward
  ML comparison, and a paper-trading loop are exactly the kind of work sample an
  energy or commodity trading desk wants to see from a candidate, independent of
  whether the included strategies are profitable.
  \item \textbf{A base for the real next step.} A power-market signal (Chapter~9)
  built on this same infrastructure -- honest costs, honest validation, no
  look-ahead -- is a genuinely useful research tool, because power markets have
  structural, physically-grounded drivers (weather, load, plant outages, storage
  economics) that are a fundamentally different estimation problem from ``predict
  next week's equity return from its own price history.''
  \item \textbf{Market literacy that compounds.} Building the cost model, the
  vol-targeting formula, and the walk-forward splitter is a faster way to
  internalise how markets and risk actually work than reading about it.
\end{enumerate}

\begin{caveat}
If a future version of this pipeline, on real data, produces a consistently
positive out-of-sample Sharpe after realistic costs: be more suspicious, not
less. Check for a data bug, a leak, or a benchmark that was quietly easy before
concluding the strategy is real.
\end{caveat}
